\documentclass[12pt,a4paper]{article}

% ========== Pacotes ==========
\usepackage[utf8]{inputenc}
\usepackage[T1]{fontenc}
\usepackage[brazil]{babel}
\usepackage{amsmath, amssymb}
\usepackage{geometry}
\usepackage{graphicx}
\usepackage{hyperref}
\usepackage{listings}
\usepackage{xcolor}
\usepackage{booktabs}
\usepackage{enumitem}
\usepackage{fancyhdr}
\usepackage{titlesec}
\usepackage{tcolorbox}

\geometry{margin=2.5cm}

% ========== Cores ==========
\definecolor{codebg}{HTML}{0F172A}
\definecolor{codetext}{HTML}{E2E8F0}
\definecolor{teal600}{HTML}{0D9488}
\definecolor{keyword}{HTML}{14B8A6}
\definecolor{string}{HTML}{34D399}
\definecolor{comment}{HTML}{64748B}

% ========== Configuração de Listings ==========
\lstset{
  basicstyle=\ttfamily\small\color{codetext},
  backgroundcolor=\color{codebg},
  keywordstyle=\color{keyword}\bfseries,
  stringstyle=\color{string},
  commentstyle=\color{comment}\itshape,
  breaklines=true,
  frame=single,
  rulecolor=\color{gray!40},
  numbers=left,
  numberstyle=\tiny\color{gray},
  numbersep=8pt,
  tabsize=2,
  showstringspaces=false,
  xleftmargin=1.5em,
  framexleftmargin=1.5em,
}

\lstdefinelanguage{json}{
  morestring=[b]",
  literate=
    *{:}{{{\color{keyword}{:}}}}{1}
    {,}{{{\color{keyword}{,}}}}{1}
    {\{}{{{\color{keyword}{\{}}}}{1}
    {\}}{{{\color{keyword}{\}}}}}{1}
    {[}{{{\color{keyword}{[}}}}{1}
    {]}{{{\color{keyword}{]}}}}{1},
}

% ========== Cabeçalho/Rodapé ==========
\pagestyle{fancy}
\fancyhf{}
\fancyhead[L]{\small Calculadora Financeira --- Documentação Técnica}
\fancyhead[R]{\small\thepage}
\fancyfoot[C]{\small\textit{Gerado automaticamente}}
\renewcommand{\headrulewidth}{0.4pt}

% ========== Capa ==========
\title{
  \vspace{2cm}
  {\Huge\bfseries Calculadora Financeira}\\[0.5cm]
  {\Large Documentação Técnica do Sistema}\\[1cm]
  {\normalsize React + Kotlin/Spring Boot}
}
\author{Projeto Acadêmico --- Matemática Financeira}
\date{Março de 2026}

% ========================================================================
\begin{document}

\maketitle
\thispagestyle{empty}
\newpage

\tableofcontents
\newpage

% ========================================================================
\section{Visão Geral do Sistema}

A \textbf{Calculadora Financeira} é uma aplicação web completa para cálculos de matemática financeira. O sistema é dividido em duas camadas:

\begin{itemize}[leftmargin=2em]
  \item \textbf{Frontend}: aplicação SPA (Single Page Application) construída com \textbf{React} + \textbf{Vite}, estilizada com \textbf{Tailwind CSS}, servida na porta \texttt{5173}.
  \item \textbf{Backend}: API REST construída com \textbf{Kotlin} + \textbf{Spring Boot 3.5}, compilada com \textbf{Gradle}, servida na porta \texttt{8080}.
\end{itemize}

\subsection{Módulos Implementados}

\begin{enumerate}
  \item \textbf{Capitalização Simples} --- cálculo de VP, VF, J, $i$ e $n$ com juros simples.
  \item \textbf{Capitalização Composta} --- cálculo de VP, VF, J, $i$ e $n$ com juros compostos.
  \item \textbf{Desconto Comercial} --- cálculo de taxa de desconto ($d$) e taxa efetiva ($i_e$).
  \item \textbf{Taxas Equivalentes} --- conversão de taxa efetiva entre períodos (diário, mensal, bimestral, trimestral, semestral, anual).
  \item \textbf{Taxas Abaixo (Nominal $\leftrightarrow$ Proporcional)} --- conversão entre taxa nominal $i_k$ e taxa proporcional $i_{\text{prop}}$.
\end{enumerate}

% ========================================================================
\section{Arquitetura}

\subsection{Padrão MVC}

Tanto o backend quanto o frontend seguem o padrão \textbf{Model--View--Controller}:

\begin{center}
\begin{tabular}{lll}
  \toprule
  \textbf{Camada} & \textbf{Backend (Kotlin)} & \textbf{Frontend (React)} \\
  \midrule
  Model      & \texttt{model/*.kt} (data classes)    & Estado local (\texttt{useState}) \\
  Service    & \texttt{service/*.kt} (lógica)         & \texttt{services/*.js} (HTTP) \\
  Controller & \texttt{controller/*.kt} (REST)        & \texttt{controllers/*.js} (parsers) \\
  View       & ---                                    & \texttt{components/*.jsx} (UI) \\
  \bottomrule
\end{tabular}
\end{center}

\subsection{Comunicação}

\begin{tcolorbox}[colback=codebg,coltext=codetext,colframe=gray!40,title=Fluxo de uma requisição]
\ttfamily\small
Componente React $\xrightarrow{\text{onSubmit}}$ Controller (JS)
$\xrightarrow{\text{fetch POST}}$ Controller (Kotlin)
$\xrightarrow{}$ Service (Kotlin) $\xrightarrow{\text{JSON}}$ Frontend
\end{tcolorbox}

O client HTTP (\texttt{httpClient.js}) encapsula \texttt{fetch} e retorna diretamente o payload JSON (sem \texttt{.data}).

\subsection{Endpoints da API}

\begin{center}
\begin{tabular}{llp{6.5cm}}
  \toprule
  \textbf{Método} & \textbf{Rota} & \textbf{Descrição} \\
  \midrule
  POST & \texttt{/api/simples/calculate}            & Capitalização simples \\
  POST & \texttt{/api/composta/calculate}            & Capitalização composta \\
  POST & \texttt{/api/desconto-comercial/calculate}  & Desconto comercial \\
  POST & \texttt{/api/taxa-equivalente/calculate}     & Taxas equivalentes \\
  POST & \texttt{/api/taxas-abaixo/calculate}         & Taxas abaixo (nominal/proporcional) \\
  \bottomrule
\end{tabular}
\end{center}

Todos os endpoints recebem e retornam \texttt{application/json}.

% ========================================================================
\section{Capitalização Simples}
\label{sec:simples}

\subsection{Modelo de Dados}

\begin{center}
\begin{tabular}{llp{6.5cm}}
  \toprule
  \textbf{Campo} & \textbf{Tipo} & \textbf{Descrição} \\
  \midrule
  \texttt{target}   & \texttt{String?} & Variável a calcular (\texttt{montante}, \texttt{capital}, \texttt{juros}, \texttt{taxa}, \texttt{tempo}) \\
  \texttt{capital}  & \texttt{Double?} & Valor Presente (VP) \\
  \texttt{montante} & \texttt{Double?} & Valor Futuro (VF) \\
  \texttt{juros}    & \texttt{Double?} & Juros (J) \\
  \texttt{taxa}     & \texttt{Double?} & Taxa em \% \\
  \texttt{tempo}    & \texttt{Double?} & Número de períodos ($n$) \\
  \bottomrule
\end{tabular}
\end{center}

\subsection{Fórmulas}

Na capitalização simples, os juros são proporcionais ao capital, à taxa e ao tempo.

\begin{tcolorbox}[colback=white,colframe=teal600,title=Fórmulas --- Juros Simples]

\textbf{Montante (Valor Futuro):}
\begin{equation}
  VF = VP \cdot (1 + i \cdot n)
\end{equation}

\textbf{Capital (Valor Presente):}
\begin{equation}
  VP = \frac{VF}{1 + i \cdot n}
\end{equation}

\textbf{Juros:}
\begin{equation}
  J = VP \cdot i \cdot n \quad \text{ou} \quad J = VF - VP
\end{equation}

\textbf{Taxa:}
\begin{equation}
  i = \frac{J}{VP \cdot n}
\end{equation}

\textbf{Tempo:}
\begin{equation}
  n = \frac{J}{VP \cdot i}
\end{equation}

\end{tcolorbox}

\subsection{Lógica de Cálculo}

O serviço \texttt{SimplesService} opera em dois modos:

\begin{enumerate}
  \item \textbf{Com \texttt{target} definido}: calcula a variável indicada. Se os campos necessários estiverem ausentes, retorna erro HTTP 400 com mensagem descritiva.
  \item \textbf{Sem \texttt{target}}: infere qual variável está faltando e a calcula automaticamente.
\end{enumerate}

\subsubsection{Exemplo de Requisição}

\begin{lstlisting}[language=json,title={POST /api/simples/calculate}]
{
  "target": "montante",
  "capital": 1000,
  "taxa": 2,
  "tempo": 10
}
\end{lstlisting}

\begin{lstlisting}[language=json,title={Resposta}]
{
  "montante": 1200.0
}
\end{lstlisting}

\textbf{Cálculo}: $VF = 1000 \times (1 + 0{,}02 \times 10) = 1000 \times 1{,}2 = 1200{,}00$.

% ========================================================================
\section{Capitalização Composta}
\label{sec:composta}

\subsection{Modelo de Dados}

Idêntico ao da capitalização simples (mesmos campos: \texttt{target}, \texttt{capital}, \texttt{montante}, \texttt{juros}, \texttt{taxa}, \texttt{tempo}).

\subsection{Fórmulas}

Na capitalização composta, os juros incidem sobre o montante acumulado (juros sobre juros).

\begin{tcolorbox}[colback=white,colframe=teal600,title=Fórmulas --- Juros Compostos]

\textbf{Montante (Valor Futuro):}
\begin{equation}
  VF = VP \cdot (1 + i)^n
\end{equation}

\textbf{Capital (Valor Presente):}
\begin{equation}
  VP = \frac{VF}{(1 + i)^n}
\end{equation}

\textbf{Juros:}
\begin{equation}
  J = VP \cdot \left[(1 + i)^n - 1\right] \quad \text{ou} \quad J = VF - VP
\end{equation}

\textbf{Taxa:}
\begin{equation}
  i = \left(\frac{VF}{VP}\right)^{\!\frac{1}{n}} - 1
\end{equation}

\textbf{Tempo:}
\begin{equation}
  n = \frac{\ln(VF / VP)}{\ln(1 + i)}
\end{equation}

\end{tcolorbox}

\subsection{Lógica de Cálculo}

O serviço \texttt{CompostaService} utiliza as funções \texttt{pow()} e \texttt{ln()} do Kotlin para exponenciação e logaritmo natural. Possui validação de domínio:
\begin{itemize}
  \item Para cálculo de taxa: $VP > 0$ e $n \neq 0$.
  \item Para cálculo de tempo: $VP > 0$, $VF > 0$ e $\ln(1 + i) \neq 0$.
\end{itemize}

\subsubsection{Exemplo de Requisição}

\begin{lstlisting}[language=json,title={POST /api/composta/calculate}]
{
  "target": "montante",
  "capital": 1000,
  "taxa": 2,
  "tempo": 12
}
\end{lstlisting}

\begin{lstlisting}[language=json,title={Resposta}]
{
  "montante": 1268.2417945625455
}
\end{lstlisting}

\textbf{Cálculo}: $VF = 1000 \times (1 + 0{,}02)^{12} = 1000 \times 1{,}02^{12} \approx 1268{,}24$.

% ========================================================================
\section{Desconto Comercial (Bancário)}
\label{sec:desconto}

\subsection{Modelo de Dados}

\begin{center}
\begin{tabular}{llp{6.5cm}}
  \toprule
  \textbf{Campo} & \textbf{Tipo} & \textbf{Descrição} \\
  \midrule
  \texttt{nominal}       & \texttt{Double?} & Valor nominal do título ($N$) \\
  \texttt{valorPresente} & \texttt{Double?} & Valor presente ($VP$) \\
  \texttt{desconto}      & \texttt{Double?} & Desconto absoluto ($D$) \\
  \texttt{taxaDesconto}  & \texttt{Double?} & Taxa de desconto ($d$) em \% \\
  \texttt{tempo}         & \texttt{Double?} & Prazo ($n$) \\
  \texttt{taxaEfetiva}   & \texttt{Double?} & Taxa efetiva ($i_e$) em \% \\
  \bottomrule
\end{tabular}
\end{center}

\subsection{Fórmulas}

\begin{tcolorbox}[colback=white,colframe=teal600,title=Fórmulas --- Desconto Comercial]

\textbf{Taxa de desconto:}
\begin{equation}
  d = \frac{D}{N \cdot n} \quad \text{ou} \quad d = \frac{1 - VP/N}{n}
\end{equation}

\textbf{Taxa efetiva equivalente:}
\begin{equation}
  i_e = \frac{d}{1 - d \cdot n}
\end{equation}

\textbf{Relação inversa:}
\begin{equation}
  d = \frac{i_e}{1 + i_e \cdot n}
\end{equation}

\end{tcolorbox}

\subsection{Lógica de Cálculo}

O serviço \texttt{DescontoComercialService} determina automaticamente o que calcular com base nos campos preenchidos:

\begin{itemize}
  \item Se faltam $d$ e $i_e$: calcula $d$ a partir de $N$, $D$ e $n$ (ou $N$, $VP$ e $n$).
  \item Se falta $i_e$ e existe $d$: calcula $i_e = d / (1 - d \cdot n)$.
  \item Se falta $d$ e existe $i_e$: calcula $d = i_e / (1 + i_e \cdot n)$.
\end{itemize}

% ========================================================================
\section{Taxas Equivalentes}
\label{sec:equivalente}

\subsection{Conceito}

Duas taxas são \textbf{equivalentes} quando aplicadas ao mesmo capital, pelo mesmo prazo total, produzem o mesmo montante em regime de juros \textbf{compostos}.

\subsection{Modelo de Dados}

\begin{center}
\begin{tabular}{llp{6.5cm}}
  \toprule
  \textbf{Campo} & \textbf{Tipo} & \textbf{Descrição} \\
  \midrule
  \texttt{taxaInformada}     & \texttt{Double?} & Taxa efetiva informada em \% \\
  \texttt{periodoInformado}  & \texttt{String?} & Período da taxa informada \\
  \texttt{periodoSolicitado} & \texttt{String?} & Período desejado para conversão \\
  \texttt{n}                 & \texttt{Double?} & Número de períodos (opcional) \\
  \bottomrule
\end{tabular}
\end{center}

\subsection{Períodos Suportados}

O sistema utiliza a \textbf{convenção comercial de 360 dias} e mapeia cada período ao seu número de ocorrências anuais:

\begin{center}
\begin{tabular}{lc}
  \toprule
  \textbf{Período} & \textbf{Ocorrências por ano} \\
  \midrule
  Anual      & 1 \\
  Semestral  & 2 \\
  Trimestral & 4 \\
  Bimestral  & 6 \\
  Mensal     & 12 \\
  Diário     & 360 \\
  \bottomrule
\end{tabular}
\end{center}

\subsection{Fórmula}

\begin{tcolorbox}[colback=white,colframe=teal600,title=Fórmula --- Taxa Equivalente]

Dado que $i$ é a taxa efetiva no período informado com $n_1$ ocorrências anuais, e deseja-se a taxa equivalente para um período com $n_2$ ocorrências anuais:

\begin{equation}
  i_{\text{eq}} = (1 + i)^{\,n_1 / n_2} - 1
\end{equation}

\textbf{Exemplo}: 2\% a.m. para taxa anual:
\[
  i_{\text{anual}} = (1 + 0{,}02)^{12/1} - 1 = 1{,}02^{12} - 1 \approx 0{,}2682 = 26{,}82\%
\]

\end{tcolorbox}

\subsection{Resposta da API}

O endpoint retorna:
\begin{itemize}
  \item \texttt{taxaEquivalente}: taxa equivalente em \% (arredondada a 6 casas).
  \item \texttt{taxaInformadaDecimal}: taxa informada em forma decimal.
  \item \texttt{taxaEquivalenteDecimal}: taxa equivalente em forma decimal.
  \item \texttt{conversao}: texto descritivo da conversão (ex: ``Taxa de 2\% ao mensal = 26.824179\% ao anual'').
\end{itemize}

% ========================================================================
\section{Taxas Abaixo (Nominal $\leftrightarrow$ Proporcional)}
\label{sec:abaixo}

\subsection{Conceito}

Na capitalização composta, a \textbf{taxa nominal} ($i_k$) é uma taxa enunciada que deve ser dividida pelo número de capitalizações ($k$) para obter a \textbf{taxa proporcional/efetiva por período} ($i_{\text{prop}}$).

\subsection{Modelo de Dados}

\begin{center}
\begin{tabular}{llp{6.5cm}}
  \toprule
  \textbf{Campo} & \textbf{Tipo} & \textbf{Descrição} \\
  \midrule
  \texttt{taxaNominal}       & \texttt{Double?} & Taxa nominal $i_k$ em \% \\
  \texttt{taxaProporcional}  & \texttt{Double?} & Taxa proporcional em \% \\
  \texttt{k}                 & \texttt{Int?}    & Periodicidade (nº de capitalizações) \\
  \bottomrule
\end{tabular}
\end{center}

\subsection{Fórmulas}

\begin{tcolorbox}[colback=white,colframe=teal600,title=Fórmulas --- Taxas Abaixo]

\textbf{Taxa proporcional a partir da nominal:}
\begin{equation}
  i_{\text{prop}} = \frac{i_k}{k}
\end{equation}

\textbf{Taxa nominal a partir da proporcional:}
\begin{equation}
  i_k = k \cdot i_{\text{prop}}
\end{equation}

\textbf{Exemplo}: Taxa nominal de 12\% a.a. com capitalização mensal ($k = 12$):
\[
  i_{\text{prop}} = \frac{12\%}{12} = 1\% \text{ ao mês}
\]

\end{tcolorbox}

% ========================================================================
\section{Stack Tecnológico}

\subsection{Backend}

\begin{center}
\begin{tabular}{ll}
  \toprule
  \textbf{Tecnologia} & \textbf{Versão/Detalhes} \\
  \midrule
  Linguagem      & Kotlin 1.9.25 \\
  Framework      & Spring Boot 3.5.7 \\
  Build Tool     & Gradle (Kotlin DSL), via Gradle Wrapper \\
  Java           & OpenJDK 21 (Temurin) \\
  Serialização   & Jackson (jackson-module-kotlin) \\
  Dev Tools      & Spring Boot DevTools (hot reload) \\
  \bottomrule
\end{tabular}
\end{center}

\subsection{Frontend}

\begin{center}
\begin{tabular}{ll}
  \toprule
  \textbf{Tecnologia} & \textbf{Versão/Detalhes} \\
  \midrule
  Framework     & React (via Vite) \\
  Bundler       & Vite 7.x \\
  Estilização   & Tailwind CSS 4.x \\
  Tipografia    & Inter (UI), JetBrains Mono (valores) \\
  HTTP Client   & \texttt{fetch} nativo (wrapper customizado) \\
  \bottomrule
\end{tabular}
\end{center}

% ========================================================================
\section{Interface do Usuário}
\label{sec:ui}

A interface apresenta um design com estética de \textbf{calculadora científica}, mantendo modernidade e responsividade.

\subsection{Estrutura Visual}

\begin{enumerate}
  \item \textbf{Cabeçalho}: ícone Σ com título ``Calculadora Financeira''.
  \item \textbf{Abas de navegação}: quatro abas --- \textbf{Simples}, \textbf{Composta}, \textbf{Desconto} e \textbf{Taxas} --- permitem alternar entre os módulos.
  \item \textbf{Tela de resultado} (\texttt{ResultDisplay}): área de display no topo de cada formulário que mostra:
  \begin{itemize}
    \item Estado vazio: ``0,00'' em cinza (como display desligado).
    \item Resultado: valor principal em destaque (fonte monospace verde, tamanho grande) com valores secundários em grade abaixo.
    \item Erro: mensagem em vermelho com ícone de alerta.
  \end{itemize}
  \item \textbf{Seletor de variável}: dropdown customizado (\texttt{Select}) para escolher qual variável calcular (target).
  \item \textbf{Campos de entrada}: inputs com design escuro, fonte monospace, com seletores de unidade (\texttt{UnitSelect}) para taxa e tempo.
  \item \textbf{Botões de ação}: ``\texttt{= Calcular}'' (gradiente teal com efeito de pressão) e ``\texttt{C}'' (limpar).
\end{enumerate}

\subsection{Componentes Reutilizáveis}

\begin{center}
\begin{tabular}{lp{8cm}}
  \toprule
  \textbf{Componente} & \textbf{Função} \\
  \midrule
  \texttt{Select}           & Dropdown customizado com teclado (setas, Enter, Esc), usa \texttt{useRef} para evitar problemas de dependência. \\
  \texttt{UnitSelect}       & Seletor de unidade (á.m., a.d., a.a., etc.) para campos de taxa e tempo. \\
  \texttt{FormulaTooltip}   & Tooltip que exibe a fórmula utilizada ao passar o mouse. \\
  \texttt{ResultDisplay}    & Display LCD-like com formatação automática de valores numéricos e percentuais. \\
  \bottomrule
\end{tabular}
\end{center}

\subsection{Conversão de Unidades no Frontend}

O frontend realiza conversão de unidades de taxa e tempo \textbf{antes} de enviar ao backend, utilizando o utilitário \texttt{timeUnits.js} com a função \texttt{convertPeriods()}. Períodos suportados: dia, mês, trimestre, semestre e ano.

% ========================================================================
\section{Tratamento de Erros}

O sistema implementa validação em ambas as camadas:

\subsection{Backend}

Os serviços de capitalização simples e composta validam as entradas e lançam \texttt{ResponseStatusException} com código HTTP \textbf{400 (Bad Request)} e mensagem descritiva quando os campos necessários estão ausentes.

\begin{lstlisting}[language=java,title={Exemplo de validação (SimplesService.kt)}]
"montante" -> when {
    vp != null && i != null && n != null ->
        mapOf("montante" to vp * (1 + i * n))
    else -> throw ResponseStatusException(
        HttpStatus.BAD_REQUEST,
        "Para calcular Montante (VF), informe
         (Capital, Taxa e Tempo) ou (Capital e Juros)."
    )
}
\end{lstlisting}

\subsection{Frontend}

O frontend captura erros da API e os exibe no componente \texttt{ResultDisplay}:

\begin{lstlisting}[language=java,title={Captura de erro no componente}]
catch (err) {
  setError(err?.payload?.message
    || err?.message
    || 'Erro ao calcular');
}
\end{lstlisting}

% ========================================================================
\section{Como Executar}

\subsection{Pré-requisitos}

\begin{itemize}
  \item Java 17+ (recomendado: OpenJDK 21)
  \item Node.js LTS + npm
\end{itemize}

\subsection{Backend (porta 8080)}

\begin{lstlisting}[language=bash,title={Terminal na pasta backend/}]
.\gradlew.bat bootRun
\end{lstlisting}

Ou via \textbf{IntelliJ IDEA}: abra a pasta \texttt{backend/}, aguarde a indexação do Gradle e clique no ícone \textbf{Run} ($\blacktriangleright$) ao lado do \texttt{fun main()} em \texttt{BackendApplication.kt}.

\subsection{Frontend (porta 5173)}

\begin{lstlisting}[language=bash,title={Terminal na pasta app/}]
npm install
npm run dev
\end{lstlisting}

Acesse \url{http://localhost:5173} no navegador.

% ========================================================================
\section{Estrutura de Diretórios}

\begin{lstlisting}[title={Árvore do projeto}]
calc-financeira/
  app/                          # Frontend React
    src/
      components/
        SimplesForm.jsx         # Formulario cap. simples
        CompostaForm.jsx        # Formulario cap. composta
        DescontoComercialForm.jsx
        TaxaEquivalenteForm.jsx
        ui/
          Select.jsx            # Dropdown customizado
          UnitSelect.jsx        # Seletor de unidade
          FormulaTooltip.jsx    # Tooltip de formula
          ResultDisplay.jsx     # Display de resultado
      controllers/
        simplesController.js
        compostaController.js
        descontoController.js
        taxaEquivalenteController.js
        taxasAbaixoController.js
      services/
        simplesService.js
        compostaService.js
        descontoService.js
        taxaEquivalenteService.js
        taxasAbaixoService.js
        httpClient.js           # Wrapper fetch
      utils/
        timeUnits.js            # Conversao de periodos
      App.jsx                   # Layout principal
      App.css                   # Estilos calculadora
      index.css                 # Tailwind + reset
      main.jsx                  # Entry point
  backend/                      # Backend Kotlin/Spring
    src/main/kotlin/com/calculadora/backend/
      BackendApplication.kt     # Classe principal
      model/
        CapitalizacaoSimples.kt
        CapitalizacaoComposta.kt
        DescontoComercial.kt
        TaxaEquivalente.kt
        TaxasAbaixo.kt
      service/
        SimplesService.kt
        CompostaService.kt
        DescontoComercialService.kt
        TaxaEquivalenteService.kt
        TaxasAbaixoService.kt
      controller/
        SimplesController.kt
        CompostaController.kt
        DescontoComercialController.kt
        TaxaEquivalenteController.kt
        TaxasAbaixoController.kt
    build.gradle.kts            # Build Gradle
    gradlew.bat                 # Gradle Wrapper (Windows)
\end{lstlisting}

% ========================================================================
\section{Resumo das Fórmulas}

\begin{center}
\renewcommand{\arraystretch}{1.6}
\begin{tabular}{lll}
  \toprule
  \textbf{Módulo} & \textbf{Fórmula Principal} & \textbf{Regime} \\
  \midrule
  Cap.\ Simples   & $VF = VP(1 + i \cdot n)$ & Simples \\
  Cap.\ Composta  & $VF = VP(1 + i)^n$       & Composto \\
  Desconto Com.   & $i_e = \dfrac{d}{1 - d \cdot n}$ & Simples \\[0.4em]
  Taxa Equiv.     & $i_{\text{eq}} = (1 + i)^{n_1/n_2} - 1$ & Composto \\[0.2em]
  Taxa Abaixo     & $i_{\text{prop}} = i_k / k$ & Proporcional \\
  \bottomrule
\end{tabular}
\end{center}

\end{document}
